% ====================================================================
\section{Synchronisation and clocks}\label{sec:clocks}
% ====================================================================

This section answers three questions about the control words. Which words can a single autonomous
tree produce (\cref{sec:phasewords,sec:generators})? How do we certify that every control word of the
processor is such a word (\cref{sec:phasecalc})? And how are the generating trees written in
logarithmic space (\cref{sec:enum})?

\subsection{Phase words}\label{sec:phasewords}

\begin{definition}[Phase words]\label{def:phaseword}
Let $\beta\in\{0,1\}$ and $\Hz\in\mathbb N$. A word $w:[0,\Hz]\to\{0,1\}$ is \emph{$\beta$-phased} if
\begin{equation}\label{eq:phase}
w(s)\ne w(s-1)\ \Longrightarrow\ w(s)=\beta\oplus(s\bmod 2)\qquad(1\le s\le\Hz).
\end{equation}
Its \emph{defect set} is $D(w,\beta)=\{t\in[0,\Hz]:\ w(t)\ne\beta\oplus(t\bmod 2)\}$.
\end{definition}

So with $\beta=0$ rises occur at odd times and falls at even times; with $\beta=1$ the parities
are reversed. Constant intervals of any length are allowed, provided that the changes at their ends
have the right parity. The initial bit $w(0)$ is unconstrained.

\begin{lemma}[Characterisation]\label{lem:charwords}
$w$ is $\beta$-phased if and only if $D(w,\beta)$ contains no two consecutive integers.
\end{lemma}

\begin{proof}
Put $\delta(t)=w(t)\oplus\beta\oplus(t\bmod 2)$, the indicator of $D$. Then
$\delta(s)\oplus\delta(s-1)=w(s)\oplus w(s-1)\oplus 1$. If $w$ retains its value at $s$, then
$\delta(s)\ne\delta(s-1)$, so $s-1,s$ are not both defects. If $w$ changes at $s$, then
$\delta(s)=\delta(s-1)$, and both are defects exactly when $w(s)\ne\beta\oplus(s\bmod 2)$. Thus a
consecutive pair of defects at $(s-1,s)$ occurs exactly when \eqref{eq:phase} fails at $s$.
\end{proof}

\begin{remark}\label{rem:notarbitrary}
Phase words are \emph{not} arbitrary synchronous words. Prescribing the samples at one parity does
not justify ignoring the updates in between. In what follows every physical operation is checked at
every transition, and every control slot keeps a single fixed phase during the whole program.
\end{remark}

\subsection{The phase calculus}\label{sec:phasecalc}

A \emph{phase assignment} gives each control slot a phase in $\{0,1\}$, such that the four slots of
a cell $X$ share one phase $e_X$. We write $e_g$ for the phase of the control of a junction $g$, and
$e_\kappa$ for the phase of an external port $\kappa$ when it is realised by a generator.

\begin{lemma}[Cell phase]\label{lem:cellphase}
Let $X$ have phase $e$ and mode word $\mu$ on $[0,\Hz]$. The four slot words of $X$ are $e$-phased if
and only if, for every $s\in[1,\Hz]$ with $\mu(s)\ne\mu(s-1)$:
\begin{enumerate}[label=(\alph*)]
\item if $\mu(s)=\FAV z$, then $s\equiv e\oplus z\pmod 2$;
\item if $\mu(s-1)=\FAV z$, then $s\equiv e\oplus z\oplus 1\pmod 2$.
\end{enumerate}
Equivalently, every maximal $\FAV z$-block that starts in $[1,\Hz]$ starts at a time of parity
$e\oplus z$; and every maximal block that ends inside the horizon, followed by a different mode, has
odd length.
\end{lemma}

\begin{proof}
The first slot word is $\one\{\mu=\LOAD\}$ and the second is $\one\{\mu\ne\RESET\}$. There are six
kinds of mode changes at $s$:
\begin{center}\small
\begin{tabular}{lll}\toprule
change at $s$ & slot changes & parity required by \eqref{eq:phase}\\\midrule
$\HOLD\to\LOAD$ & first slot rises & $s\equiv e\oplus1$\\
$\LOAD\to\HOLD$ & first slot falls & $s\equiv e$\\
$\HOLD\to\RESET$ & second slot falls & $s\equiv e$\\
$\RESET\to\HOLD$ & second slot rises & $s\equiv e\oplus 1$\\
$\RESET\to\LOAD$ & both rise & $s\equiv e\oplus1$\\
$\LOAD\to\RESET$ & both fall & $s\equiv e$\\\bottomrule
\end{tabular}
\end{center}
In every row the required parity is the one demanded by (a) for the new mode, if it is not \HOLD,
and by (b) for the old mode, if it is not \HOLD. When both apply, as in the last two rows, they demand
the same parity. Conversely (a) and (b) give exactly the parities of the table. A block that starts
at parity $e\oplus z$ and ends at parity $e\oplus z\oplus 1$ has odd length.
\end{proof}

\begin{lemma}[Single-bit controls]\label{lem:bitphase}
A junction control or an external port of phase $e$ has an $e$-phased word if and only if each change
of its value to $v$ happens at a time $s\equiv e\oplus v$.
\end{lemma}

\begin{proof}
This is \eqref{eq:phase} restated.
\end{proof}

\begin{definition}[Blocks and admissibility]\label{def:block}
A \emph{block} is a triple $(X,[s,s+\ell),z)$ meaning that cell $X$ is in $\FAV z$ during
$[s,s+\ell)$. It is \emph{admissible} if $\ell$ is odd and $s\equiv e_X\oplus z$. A \emph{control
event} $(g,s,v)$ means that the control of $g$ changes to $v$ at time $s$. It is \emph{admissible} if
$s\equiv e_g\oplus v$, and similarly for external ports.
\end{definition}

\begin{lemma}[Concatenation]\label{lem:concat}
Suppose that the mode word of a cell $X$ on $[0,\Hz]$ is \HOLD\ outside a family of pairwise
disjoint admissible blocks, all contained in $[0,\Hz)$. Then the slot words of $X$ are
$e_X$-phased. In particular, two admissible blocks of the same mode are never adjacent, so blocks
never merge. (A block starting at time $0$ causes no change at time $0$; its start parity is
admissible anyway, and only its end is constrained by \cref{lem:cellphase}.)
\end{lemma}

\begin{proof}
Let $(X,[s,s+\ell),z)$ and $(X,[s+\ell,s+\ell+\ell'),z')$ be adjacent admissible blocks. Then
$s+\ell\equiv e_X\oplus z\oplus 1$ because $\ell$ is odd, and $s+\ell\equiv e_X\oplus z'$ by
admissibility. Hence $z'\ne z$. So the maximal constant-mode intervals that are not \HOLD\ are
exactly the given blocks. Each starts at parity $e_X\oplus z$ and has odd length, and
\cref{lem:cellphase} applies.
\end{proof}

\begin{remark}\label{rem:merging}
Without the concatenation lemma, a subtle failure would be possible. Two three-step $\FAV z$ phases
scheduled back to back on the same cell would form one block of length six, whose end has the wrong
parity. The original manuscript does not discuss this case. \Cref{lem:concat} shows that it cannot
arise once every scheduled block is admissible.
\end{remark}

\paragraph{Timing signatures.}
A primitive is a schedule fragment given relative to its start $s$: a list of blocks with offsets
$(X,[s+a,s+a+\ell),z)$ and of control events $(g,s+a,v)$. In background $c$ every mode and every
control value is of the form $z=c\oplus\zeta$ with $\zeta\in\{0,1\}$; here $\zeta=0$ means
``favour/assert the background''. Admissibility of an event of phase $e$ at offset $a$ then reads
\begin{equation}\label{eq:startparity}
s\equiv e\oplus c\oplus\zeta\oplus a\pmod 2 .
\end{equation}
A primitive has \emph{reference parity} $\pi$ if all its events are admissible exactly when
$s\equiv\pi$. By \eqref{eq:startparity}, the reference parity of the background-$c$ instance is that
of the background-$0$ instance shifted by $c$. This is how duality (\cref{lem:duality}) acts on
timings.

\begin{lemma}[Timing signatures]\label{lem:timing}
With the local phase conventions of \cref{tab:timing}, each primitive has the stated duration and
reference parity, and all its blocks have odd length.
\end{lemma}

\begin{table}[ht]
\centering\small
\begin{tabular}{>{\raggedright\arraybackslash}p{4.2cm}>{\raggedright\arraybackslash}p{4.3cm}cc}\toprule
primitive & local phases & duration & reference parity\\\midrule
$\prim{SET}(X\mathbin{:=}c\oplus\zeta)$ & $e_X$ & $3$ & $e_X\oplus c\oplus\zeta$\\
$\prim{COPY}$ into $X$, mode $\FAV{c\oplus\zeta}$ & $e_X$ & $3$ & $e_X\oplus c\oplus\zeta$\\
$\prim{MOVE}_{c\oplus\zeta}(X\to Y)$ & $e_Y=e_X\oplus 1$ & $3$ & $e_X\oplus c\oplus\zeta$\\
$\prim{ROUTE}(g;A\to B;C)$ & arms $e$, $h_g$: $e\oplus1$ & $6$ & $e\oplus c$\\
$\prim{GATE}_c(g;A,B\to C)$ & arms $e$, $h_g$: $e\oplus1$ & $5$ & $e\oplus c\oplus 1$\\
$\prim{READ}(g;A;C,K,R)$ & $A,B,C,R$: $e$; $K,h_g$: $e\oplus1$ & $14$ & $e\oplus c$\\\bottomrule
\end{tabular}
\caption{Timing signatures. The reference parity is the parity of the start time $s$ for which all
blocks and control events of the primitive are admissible.}
\label{tab:timing}
\end{table}

\begin{proof}
We list each event as (phase $e'$, $\zeta$, offset $a$) and evaluate $e'\oplus c\oplus\zeta\oplus a$
by \eqref{eq:startparity}.
\begin{itemize}
\item \prim{SET}, \prim{COPY}: one block $(e_X,\zeta,0)$ of length $3$.
\item \prim{MOVE}: $X$ gives $(e_X,\zeta,0)$; $Y$ in $\FAV{1-z}$ gives
  $(e_X\oplus1,\zeta\oplus 1,0)$. Both give $e_X\oplus c\oplus\zeta$. Lengths are $3$.
\item \prim{ROUTE}: source $(e,0,0)$; destination $(e,1,3)$; $h_g\to b$ is $(e\oplus1,1,0)$; $h_g\to c$
  is $(e\oplus1,0,3)$. All four give $e\oplus c$. Lengths are $3$.
\item \prim{GATE}: output $(e,1,0)$, of length $5$; operands $(e,0,1)$, of length $3$; $h_g\to b$ is
  $(e\oplus1,1,1)$; $h_g\to c$ is $(e\oplus1,0,4)$. All give $e\oplus c\oplus1$.
\item \prim{READ}: $h_g\to b$ is $(e\oplus1,1,0)$ and $h_g\to c$ is $(e\oplus1,0,11)$. $K$ has blocks
  $(e\oplus1,0,1)$ and $(e\oplus1,1,4)$; $C$ has $(e,0,4)$ and $(e,1,11)$; $A$ has $(e,0,8)$ and
  $(e,1,11)$. All give $e\oplus c$. All blocks have length $3$, and the two blocks of $K$, and the
  two blocks of $A$, are adjacent with different modes.\qedhere
\end{itemize}
\end{proof}

\subsection{Robust generators: constants, timers and pulses}\label{sec:timers}

Robust generators were defined in \cref{def:robust}. All branches below are private, disjoint copies.

\paragraph{Constant sources.} A \emph{constant source} of bit $c$ consists of a root and two leaves,
all initially $c$. With its consumer the root has degree three: its two leaves always vote $c$, so
it stays $c$, and the leaves copy the root. It is a robust generator of the constant word $c$, with
three vertices and degrees $1$ and $3$.

\begin{lemma}[Timers]\label{lem:timers}
For every $h\ge 0$ there is a robust generator $D_h$ of $t\mapsto\one\{t<h\}$, with $8h-1$ vertices
if $h\ge 1$ and $7$ if $h=0$, using only degrees $1$ and $3$, consumer edge included. Complementing
all initial bits gives a robust generator $U_h$ of $\one\{t\ge h\}$.
\end{lemma}

\begin{proof}
For $h=0$, take a root initially $0$ with two constant-zero sources. For $h\ge 1$, initialise the
root to $1$ and put $L=h-1$. If $L=0$, attach two constant-zero sources to the root. If $L\ge 1$,
attach two identical chains $p_{L-1},\dots,p_0$, numbered from the bottom $p_0$, all initially
$1$. The head $p_{L-1}$ is adjacent to the root. For $k\ge 1$, $p_k$ is adjacent to its parent,
its child $p_{k-1}$ and one constant-zero source. The bottom $p_0$ is adjacent to its parent and two
constant-zero sources. Every chain vertex has degree three and computes
$\MAJ_3(\text{parent},\text{child},0)=\text{parent}\wedge\text{child}$, where the bottom's missing
child is read as $0$.

We prove by induction on $t$, simultaneously for all vertices and for an arbitrary consumer history,
that
\[
x_{p_k}(t)=\one\{t<k+1\},\qquad x_{\mathrm{root}}(t)=\one\{t<h\}.
\]
At $t=0$ all these bits are $1$. For the step: $x_{p_0}(t+1)=0=\one\{t+1<1\}$. For $k\ge 1$, the
child gives $x_{p_{k-1}}(t)=\one\{t<k\}$. The parent is $p_{k+1}$ with history $\one\{t<k+2\}$, or the
root with history $\one\{t<L+1\}$, which is the same expression when $k=L-1$. In both cases the
parent's bit is $1$ whenever $t<k$, so the conjunction equals $\one\{t<k\}=\one\{t+1<k+1\}$. The two
heads have equal histories, so the root computes $\MAJ_3(\text{head},\text{head},y(t))$, which is the
head's bit $\one\{t<L\}=\one\{t+1<h\}$ whatever $y(t)$ is. The same induction covers $L=0$.

The size is $1+2\bigl(L+3(L-1)+6\bigr)=8h-1$ for $L\ge1$, and $7$ for $h\in\{0,1\}$. Complementation
preserves majority at odd degree (\cref{lem:selfdual}), and maps arbitrary consumer histories
bijectively to arbitrary consumer histories.
\end{proof}

No robustness is claimed for a chain head on its own. The construction duplicates synchronised
chains; it does not recursively duplicate robust timers.

\begin{lemma}[Pulses]\label{lem:pulses}
For integers $0\le a<b$ there is a robust generator $P_{a,b}$ of $t\mapsto\one\{a\le t<b\}$ with
$O((a+b+1)^2)$ vertices, using only degrees $1$ and $3$. Its history is exact for all $t\ge 0$; in
particular it is $0$ for all $t\ge b$.
\end{lemma}

\begin{proof}
If $a=0$, use $D_b$. Otherwise put $B=b-a\ge 1$ and $L=a-1$, and initialise the root to $0$. If
$L=0$, its two children are private copies of $D_B$, so the root's bit at $t+1$ is $D_B(t)$, and its
history is $\one\{1\le t<B+1\}$. If $L\ge 1$, attach two identical chains $p_{L-1},\dots,p_0$,
initially $0$, with the head adjacent to the root. $p_0$ is adjacent to its parent and two copies of
$D_B$. For $k\ge1$, $p_k$ is adjacent to its parent, to $p_{k-1}$, and to one copy of $D_{B+k}$. We
claim
\[
x_{p_k}(t)=\one\{k+1\le t<B+k+1\},\qquad x_{\mathrm{root}}(t)=\one\{L+1\le t<B+L+1\}=\one\{a\le t<b\}.
\]
At $t=0$ all these bits are $0$. The bottom copies its two equal timers:
$x_{p_0}(t+1)=D_B(t)$. For $k\ge 1$ consider control time $t$:
\begin{itemize}
\item if $t<k$: the child is $0$. The parent is $p_{k+1}$, with history $\one\{k+2\le t<\dots\}$, or
  the root, with history $\one\{L+1\le t<\dots\}$, which is the same when $k=L-1$; so it is $0$ too.
  Hence $x_{p_k}(t+1)=0$;
\item if $t=k$: the child is $\one\{k\le k<B+k\}=1$ and the timer is $\one\{k<B+k\}=1$, so
  $x_{p_k}(k+1)=1$;
\item if $k<t<B+k$: the child and the timer are $1$;
\item if $t\ge B+k$: the child and the timer are $0$.
\end{itemize}
Only the first case uses the parent, through the induction hypothesis. The two heads agree, so the
root copies them whatever its consumer does. One branch has $L$ chain vertices, two copies of $D_B$,
and one $D_{B+k}$ for each $1\le k<L$. Its size is
$O\bigl(L+B+\sum_{k=1}^{L-1}(B+k)\bigr)=O((a+b+1)^2)$.
\end{proof}

\subsection{The phase-defect theorem}\label{sec:generators}

\begin{theorem}[Phase-defect generators]\label{thm:phasedefect}
Let $D\subseteq\mathbb N$ be finite with no two consecutive elements. Put $K=|D|$, fix $\Hz\ge\max D$
($\Hz=0$ if $D=\emptyset$), and let $\beta\in\{0,1\}$. There is a robust generator $\Gamma(D,\beta)$
of
\[
h(t)=\beta\oplus(t\bmod 2)\oplus\one\{t\in D\}\qquad(t\ge 0),
\]
with all degrees in $\{1,3,4\}$, consumer edge included, and size
\begin{equation}\label{eq:gensize}
|\Gamma(D,\beta)|=O\bigl((K+1)^{\alpha}(\Hz+1)^3\bigr),\qquad \alpha=\frac{\log 3}{\log(3/2)}<3 .
\end{equation}
\end{theorem}

The output alternates after the last defect: the theorem does not promise a constant tail. The
proof uses one recursive step (the combiner) and three base cases.

\begin{lemma}[Combiner]\label{lem:combiner}
Let $D$ and $\beta$ be as in \cref{thm:phasedefect}, and let $E=D\setminus\{0\}$, with its elements
ranked $0,1,2,\dots$ in increasing order. For $j\in\{0,1,2\}$ put
\begin{equation}\label{eq:childsets}
D_j=\{a-1:\ a\in E,\ \rank_E(a)\not\equiv j\pmod 3\}.
\end{equation}
Let $\gamma$ be a vertex with initial bit $\beta\oplus\one\{0\in D\}$, whose three children are the
roots of robust generators $\Gamma_j$ of $g_j(t)=(1-\beta)\oplus(t\bmod 2)\oplus\one\{t\in D_j\}$. Then
the tree rooted at $\gamma$ is a robust generator of $h$, and $\gamma$ has degree four with its
consumer. Each $D_j$ has no two consecutive elements, $\max D_j\le\max D-1$, and every element of
$E$ is passed, shifted by $-1$, to exactly two of the three children.
\end{lemma}

\begin{proof}
For $t\ge0$ we have $t+1\ge1$, so $t+1\in D$ if and only if $t+1\in E$, and
\[
\begin{aligned}
g_j(t)&=\beta\oplus\bigl((t+1)\bmod 2\bigr)\oplus\one\{t+1\in E,\ \rank(t+1)\not\equiv j\}\\
&=h(t+1)\oplus\one\{t+1\in E,\ \rank(t+1)\equiv j\}.
\end{aligned}
\]
So all three children equal $h(t+1)$ when $t+1\notin D$, and exactly two of them do when $t+1\in D$.
We prove $x_\gamma(t)=h(t)$ by induction on $t$, for an arbitrary consumer history. The base case is
the initial bit. For the step, $x_\gamma(t+1)$ is $\MAJ_5$ of $x_\gamma(t)$, the consumer and the
three children:
\begin{itemize}
\item if $t+1\notin D$, the three children vote $h(t+1)$, which decides $\MAJ_5$;
\item if $t+1\in D$, then $t\notin D$ (no consecutive defects), so
  $h(t+1)=\beta\oplus(t+1)\oplus1=\beta\oplus t=h(t)=x_\gamma(t)$. Two children and the self-vote give
  three votes for $h(t+1)$.
\end{itemize}
The children are robust and their consumer is $\gamma$, so their histories are as stated. Each $D_j$
is a subset of the shift $E-1$ of a set with no consecutive elements, which proves the remaining
claims.
\end{proof}

The original manuscript distinguishes three cases for the transition at $t$. The proof above merges
the two cases with $t+1\notin D$, because in both of them all three children agree with the target.

\begin{lemma}[Parity transformation]\label{lem:paritytransform}
Let $\Gamma$ be a robust generator of $h$ with root $\gamma$, in which every vertex has odd degree,
consumer edge included. Let $\chi(v)=\dist(v,\gamma)\bmod 2$, and let $\Gamma^\beta$ be obtained by
complementing the initial bit of every $v$ with $\chi(v)\oplus\beta=1$. Then $\Gamma^\beta$ is a robust
generator of $h(t)\oplus(t\bmod2)\oplus\beta$.
\end{lemma}

\begin{proof}
Let $M_v(t)=(t\bmod 2)\oplus\chi(v)\oplus\beta$, and given a trajectory $x$ of $\Gamma$ put
$y_v(t)=x_v(t)\oplus M_v(t)$. If $u$ is adjacent to $v$ then $\chi(u)=\chi(v)\oplus1$, so
$M_u(t)=M_v(t+1)$. At odd degree there is no self-vote, and $\MAJ$ commutes with complementing all
arguments. Hence
\[
y_v(t+1)=\MAJ\bigl(x_u(t)\bigr)_u\oplus M_v(t+1)=\MAJ\bigl(x_u(t)\oplus M_u(t)\bigr)_u=\MAJ\bigl(y_u(t)\bigr)_u,
\]
so $y$ is a trajectory of $\Gamma^\beta$, with $y(0)$ the stated initial bits. The consumer is at odd
distance, and $y_\epsilon(t)=x_\epsilon(t)\oplus(t+1+\beta)\bmod 2$ is a bijection between consumer
histories. So every consumer history of $\Gamma^\beta$ arises from one of $\Gamma$, and
$y_\gamma(t)=h(t)\oplus(t\bmod2)\oplus\beta$.
\end{proof}

The transformation is applied only to the odd-degree pulse trees, never to a degree-four combiner,
where the self-vote would carry the wrong mask.

\begin{proof}[Proof of \cref{thm:phasedefect}]
We first give the base cases $K\le 2$, then the recursion.

\emph{$K=0$.} The root has initial bit $\beta$ and two leaves with initial bit $1-\beta$. With its
consumer the root has degree three. The leaves copy the root, and the root copies the two equal
leaves, so the root alternates starting from $\beta$ whatever the consumer does. Three vertices.

\emph{$K=1$, $D=\{a\}$.} Apply \cref{lem:paritytransform} to the pulse $P_{a,a+1}$ of
\cref{lem:pulses}, which has degrees $1$ and $3$. The root history becomes
$\one\{t=a\}\oplus(t\bmod2)\oplus\beta=h(t)$. Size $O((a+1)^2)$.

\emph{$K=2$, $D=\{a,b\}$ with $b-a\ge 2$.} Apply the combiner. Since $0\notin D$ when $a>0$, the ranks
are $\rank(a)=0$ and $\rank(b)=1$, and the child sets are
\[
D_0=\{b-1\},\qquad D_1=\{a-1\},\qquad D_2=\{a-1,b-1\}.
\]
Only the third child needs another two-event construction. After $\ell\le a$ such steps, the
continuation has defect set $\{a-\ell,b-\ell\}$, phase $\beta\oplus(\ell\bmod2)$, and initial bit
$\beta\oplus(\ell\bmod 2)\oplus\one\{a=\ell\}$. At $\ell=a$ the set is $\{0,b-a\}$. Zero is removed
\emph{before} ranking, so the child sets are $\emptyset$, $\{b-a-1\}$ and $\{b-a-1\}$, all handled by
the bases $K\le1$. In total there are $a+1$ combiners and $2(a+1)$ bases with zero or one defect,
each of size $O((b+1)^2)$, so the size is $O((a+1)(b+1)^2)=O((\Hz+1)^3)$. The chain is indexed by
$\ell=0,\dots,a$, which gives an explicit enumeration without a recursion stack of depth $a$
(\cref{sec:enum}).

\emph{$K\ge3$.} Apply the combiner recursively until the sets have at most two events. Let
$\sigma=\one\{0\in D\}$ and $n=|E|=K-\sigma$. The number of ranks in $[0,n)$ congruent to $j$ is at
least $\lfloor n/3\rfloor$. Hence
\[
K_{\text{child}}\le n-\lfloor n/3\rfloor=\Bigl\lceil\frac{2(K-\sigma)}{3}\Bigr\rceil\le\frac{2K+2}{3},
\qquad K_{\text{child}}-2\le\frac23(K-2).
\]
A node at depth $d$ of the recursion has at most $2+(2/3)^d(K-2)$ events. A node is expanded by a
further combiner only if it has at least three events, so its depth is at most $\log_{3/2}(K-2)$. The
expanded nodes form a ternary tree with at most
$\sum_{d\le\log_{3/2}(K-2)}3^d\le\tfrac32(K-2)^{\alpha}$ nodes. The terminal nodes, those with at
most two events, number at most three times as many, and each costs $O((\Hz+1)^3)$, because all
defects stay in $[0,\Hz]$. This proves \eqref{eq:gensize}.

\emph{Topology and robustness.} Each component is joined to its private child trees by single
edges, so the result is a tree. Combiners have degree four; the bases with zero or one defect have
degrees $1$ and $3$. Robustness follows by induction from the bases upwards, using
\cref{lem:combiner,lem:paritytransform}.
\end{proof}

\begin{corollary}[Realised words]\label{cor:words}
Let $w:[0,\Hz]\to\{0,1\}$ be $\beta$-phased. Then $\Gamma\bigl(D(w,\beta),\beta\bigr)$ is a robust
generator whose root history equals $w(t)$ for every $t\in[0,\Hz]$, including $t=0$. It has degrees
in $\{1,3,4\}$ and size $O\bigl((\Hz+2)^{\alpha+3}\bigr)\subseteq O\bigl((\Hz+2)^6\bigr)$. After $\Hz$
its history alternates.
\end{corollary}

\begin{proof}
By \cref{lem:charwords}, $D=D(w,\beta)$ has no two consecutive elements, so \cref{thm:phasedefect}
applies. For $t\le\Hz$, $h(t)=\beta\oplus t\oplus\one\{t\in D\}=w(t)$. The event at time zero, if any,
is absorbed in the initial bit. Finally $K\le\Hz+1$.
\end{proof}

\subsection{Logspace enumeration of the generators}\label{sec:enum}

\begin{lemma}[Uniformity]\label{lem:genlogspace}
There is a deterministic procedure which, given $\Hz$, $\beta$ and access to a membership oracle for
$D\subseteq[0,\Hz]$, lists the vertices, edges and initial bits of $\Gamma(D,\beta)$ using
$O(\log(\Hz+2))$ bits of work space, not counting the oracle. Vertex names have $O(\log(\Hz+2))$
bits.
\end{lemma}

\begin{proof}
\emph{Skeleton addresses.} A node of the combiner recursion is named by its address
$p=(j_1,\dots,j_d)\in\{0,1,2\}^d$. Its defect set is $D_p$, with $D_{()}=D$ and
$D_{p\,j}=(D_p)_j$ as in \eqref{eq:childsets}. Expanded nodes have depth at most $\log_{3/2}K+1$, so
addresses have $O(\log(K+2))$ symbols.

\emph{Counting recurrence.} Let $Q_p(t)=|D_p\cap[0,t]|$, with $Q_p(t)=0$ for $t<0$. Let
$\sigma_p=Q_p(0)=\one\{0\in D_p\}$, and for $u\ge0$ let
\[
F_j(u)=u-\Bigl\lfloor\frac{u+2-j}{3}\Bigr\rfloor\qquad(j=0,1,2),
\]
which counts the ranks in $[0,u)$ that are not congruent to $j$. Then
\begin{equation}\label{eq:counting}
Q_{p\,j}(t)=F_j\bigl(Q_p(t+1)-\sigma_p\bigr)\qquad(t\ge 0).
\end{equation}
Indeed, the elements of $D_{pj}$ in $[0,t]$ are the values $a-1$ for $a\in E_p\cap[1,t+1]$ with
$\rank(a)\not\equiv j$. The set $E_p\cap[1,t+1]$ consists of the $u=Q_p(t+1)-\sigma_p$ elements of
smallest rank, namely the ranks $0,\dots,u-1$.

\emph{Evaluation.} Let $p_r$ be the prefix of $p$ of length $r$. Iterating \eqref{eq:counting},
$Q_{p}(t)$ is obtained from $u=Q(t+d)$ by applying $u\leftarrow F_{j_r}(u-\sigma_{p_{r-1}})$ for
$r=1,\dots,d$. The bits $\sigma_{p_0},\dots,\sigma_{p_{d-1}}$ are computed in this order, each by
the same procedure with $t=0$ applied to the prefix. Thus only the address, $d$ zero-event bits, and
a constant number of counters of $O(\log(\Hz+K+2))$ bits are stored. The root count $Q(t)$ is
obtained by scanning $[0,t]$ with the oracle. Starting from $u=K$ instead of $Q(t+d)$ gives $|D_p|$.
At a terminal node, the at most two dates are found by scanning $t=0,\dots,\Hz$ and comparing
$Q_p(t)$ with $Q_p(t-1)$.

\emph{Emission.} Enumerate the addresses of length at most $4\lceil\log_2(K+2)\rceil+4$, and keep
those whose proper prefixes all have more than two events. There are polynomially many candidates.
For each kept node, emit a combiner, with phase $\beta\oplus(d\bmod2)$ and initial bit
$\beta\oplus(d\bmod2)\oplus\sigma_p$, and its edges to its children; or, at a terminal node, emit the
base:
\begin{itemize}
\item a two-event base: its explicit chain indexed by $\ell$, with the phases and initial bits given
  in the proof of \cref{thm:phasedefect};
\item a one-event base: the pulse tree of \cref{lem:pulses}. Each name records branch, chain index,
  timer index and position inside the timer and constant sources, so the distance parity to the pulse
  root, and hence the complementation of \cref{lem:paritytransform}, is computed arithmetically;
\item a zero-event base: three vertices.
\end{itemize}
Every name consists of an address, a base type and $O(1)$ indices bounded by $\Hz+2$.
\end{proof}

\begin{remark}\label{rem:largedates}
For dates that are polynomially bounded in an external input size, and a logspace membership
predicate, composition gives ordinary logspace uniformity in that input size. For arbitrarily large
binary dates the lemma gives only the stated space bound in terms of $\Hz$. We do not claim
polynomial size in $\log\Hz$ for exponentially large dates.
\end{remark}
